\section*{Invariantes y cuadrivectores}
\begin{align*}
    &\text{Elemento de línea o intervalo espacio-temporal: } ds^2 \equiv g_{\alpha \beta}dx^\alpha dx^\beta\\
        &ds^2 >0 : \text{tipo espacio}\; \; \; ds^2<0:\text{tipo tiempo} \;\;\; ds^2 =0 : \text{tipo luz}\\
    &ds^2 =-c^2 d\tau ^2\equiv \text{Invariante}\; \; \;  \tau \equiv \text{tiempo propio} \;\;\; d\tau \equiv \text{Invariante}\\
    &\text{Cuadrivector velocidad: } \underline{u} = \frac{dx^i}{d\tau} e_i \; \; \; \underline{u} \cdot \underline{u} = -c^2 \;\; \; \frac{dx}{dt}=c\frac{u^1}{u^0}\; \; \; \underline{a}\cdot \underline{u} =0\\
    &\text{Cuadrivector aceleración: } \underline{a}\!=\! \frac{d\underline{u}}{d\tau}\!=\! \left(\frac{du^\alpha }{d\tau}+\Gamma^\alpha_{\beta \gamma}u^\beta u^\gamma\right) e_\alpha\;\; \; \text{Truco: }\!\frac{d\vec{A}}{d\tau}\!=\!\frac{dx^\alpha}{d\tau}\frac{\partial \vec{A}}{\partial x^\alpha}\\ 
    &\text{Aceleración propia: }a = \sqrt{\underline{a}\cdot \underline{a}}\;  \text{ (aceleración medida por observador comóvil).}\\
    &\text{Sistema comóvil: marco que se mueve a la misma velocidad que la partícula.}
\end{align*}
